\documentclass[12pt]{article}

\usepackage{fontspec}
\newfontfamily\handwriting{a-casual-handwritten-pen}[Path=../assets/fonts/,Extension=.ttf,Scale=1.5]
\newfontfamily\comic{Comic Neue}

\usepackage[utf8]{inputenc}
\usepackage{amsmath}
\usepackage{amssymb}
\usepackage{enumitem}
\usepackage{xcolor}
\usepackage{soul}
\usepackage{tikz}
\usepackage[most]{tcolorbox}
\usepackage{titlesec}
\usepackage{multicol}
\usepackage{environ}

\usetikzlibrary{fit, positioning, arrows.meta, decorations.pathreplacing}

% --- Custom Colors ---
\definecolor{primaryblue}{RGB}{42, 111, 151}
\definecolor{exbg}{RGB}{255, 240, 245}
\definecolor{rulebg}{RGB}{232, 245, 233}
\definecolor{highlight}{RGB}{255, 243, 176}
\definecolor{pinkanno}{RGB}{255, 105, 180}
\definecolor{purpleanno}{RGB}{147, 112, 219}
\definecolor{solutionbg}{RGB}{245, 245, 255}

% --- Header Styling ---
\titleformat{\section}{\color{primaryblue}\normalfont\Large\bfseries}{\thesection}{1em}{}
\titleformat{\subsection}{\color{primaryblue}\normalfont\large\bfseries}{\thesubsection}{1em}{}

% --- Friendly Boxes ---
% --- Auto-numbered Example Box ---
\newcounter{examplenum}

\newtcolorbox{friendlyexbox}[1]{
    colback=exbg,
    colframe=pinkanno!50,
    arc=5pt,
    title=\textbf{#1},
    coltitle=black,
    attach title to upper,
    after title={:\quad}
}

\newenvironment{friendlyex}{%
    \refstepcounter{examplenum}%
    \begin{friendlyexbox}{Example \theexamplenum}%
}{%
    \end{friendlyexbox}%
}

\newtcolorbox{friendlyrule}[1]{
    colback=rulebg,
    colframe=primaryblue!30,
    arc=5pt,
    fonttitle=\bfseries,
    title={#1},
    coltitle=primaryblue
}

\newcommand{\funhl}[1]{\sethlcolor{highlight}\hl{#1}}

% --- Show/Hide Solutions ---
\newif\ifshowsolutions

% Uncomment the next line to show solutions.
\showsolutionstrue

\newsavebox{\solutionmeasurebox}

\NewEnviron{solutionbox}{
\ifshowsolutions
\begin{tcolorbox}[
    colback=solutionbg,
    colframe=purpleanno!45,
    arc=5pt,
    title={Solution},
    coltitle=primaryblue,
    fonttitle=\bfseries
]
\BODY
\end{tcolorbox}
\else
\sbox{\solutionmeasurebox}{\begin{minipage}{\dimexpr\linewidth-10pt\relax}\BODY\end{minipage}}
\vspace{\dimexpr\ht\solutionmeasurebox+\dp\solutionmeasurebox+3em\relax}
\fi
}

\begin{document}
\handwriting

\begin{center}
    \Large\textbf{\color{primaryblue}Module 3 Lesson 3\\Rational Functions} \\
    \vspace{0.2cm}
    \hrule
\end{center}

\section*{Objectives}

\begin{friendlyrule}{In this lesson, we will learn how to}
\begin{itemize}
    \item find the domain of a rational function;
    \item distinguish a vertical asymptote from a removable hole;
    \item find horizontal asymptotes using the degree of the numerator and denominator;
    \item find where a graph crosses its horizontal asymptote.
\end{itemize}
\end{friendlyrule}

\section*{1. Domain of a Rational Function}

\begin{friendlyrule}{Rational Function}
A \funhl{rational function} has the form $f(x)=\dfrac{p(x)}{q(x)}$, where $p(x)$ and $q(x)$ are polynomials and $q(x)\neq 0$. The domain excludes every $x$ for which $q(x)=0$.
\end{friendlyrule}

\begin{friendlyex}{}
Find the domain of each rational function.
\begin{enumerate}[label=(\alph*)]
    \item $f(x)=\dfrac{x+3}{x^2-9}$
    \item $g(x)=\dfrac{2x}{x^2+x-6}$
\end{enumerate}
\end{friendlyex}

\begin{solutionbox}
\textbf{(a)} $x^2-9=(x-3)(x+3)=0 \implies x=3,-3$. Domain: all reals except $x=3,-3$, i.e.\ $\{x\mid x\neq \pm 3\}$.

\textbf{(b)} $x^2+x-6=(x+3)(x-2)=0 \implies x=-3,2$. Domain: all reals except $x=-3,2$.
\end{solutionbox}

\section*{2. Vertical Asymptotes vs.\ Holes}

\begin{friendlyrule}{Finding Vertical Asymptotes and Holes}
Factor the numerator and denominator completely.
\begin{itemize}
    \item If a factor $(x-c)$ \funhl{cancels} from both numerator and denominator, the graph has a \funhl{hole} at $x=c$ (a single missing point).
    \item If a factor $(x-c)$ remains in the denominator after cancelling, the graph has a \funhl{vertical asymptote} at $x=c$.
\end{itemize}
\end{friendlyrule}

\begin{friendlyex}{}
For each function, identify any holes and vertical asymptotes.
\begin{enumerate}[label=(\alph*)]
    \item $f(x)=\dfrac{x-2}{x^2-4}$
    \item $g(x)=\dfrac{x+1}{x^2-1}$
    \item $h(x)=\dfrac{x^2-4}{x^2+5x+6}$
\end{enumerate}
\end{friendlyex}

\begin{solutionbox}
\textbf{(a)} $f(x)=\dfrac{x-2}{(x-2)(x+2)}=\dfrac{1}{x+2}$ (for $x\neq 2$). \\
Hole at $x=2$; vertical asymptote at $x=-2$.

\textbf{(b)} $g(x)=\dfrac{x+1}{(x+1)(x-1)}=\dfrac{1}{x-1}$ (for $x\neq -1$). \\
Hole at $x=-1$; vertical asymptote at $x=1$.

\textbf{(c)} $h(x)=\dfrac{(x-2)(x+2)}{(x+2)(x+3)}=\dfrac{x-2}{x+3}$ (for $x\neq -2$). \\
Hole at $x=-2$; vertical asymptote at $x=-3$.
\end{solutionbox}

\section*{3. Horizontal Asymptotes}

\begin{friendlyrule}{Horizontal Asymptote Rules}
Compare the degree of the numerator $n$ to the degree of the denominator $m$.
\begin{itemize}
    \item If $n<m$: horizontal asymptote is $y=0$.
    \item If $n=m$: horizontal asymptote is $y=\dfrac{\text{leading coefficient of numerator}}{\text{leading coefficient of denominator}}$.
    \item If $n>m$: there is \textbf{no} horizontal asymptote (the function may have a slant asymptote instead).
\end{itemize}
\end{friendlyrule}

\section*{4. Reading Graphs of Rational Functions}

\begin{friendlyex}{}
The graph of $f(x)=\dfrac{3x}{x(x+2)(x-1)}$ is shown below. Identify the vertical asymptotes, the horizontal asymptote, and any holes.

\begin{center}
\begin{tikzpicture}
\begin{scope}[xscale=0.85,yscale=0.55]
\draw[dashed, gray] (-2,-4.5) -- (-2,4.5);
\draw[dashed, gray] (1,-4.5) -- (1,4.5);
\draw[dashed, gray] (-4.3,0) -- (4.3,0);
\draw[-{Latex}] (-4.3,0) -- (4.3,0) node[right] {\small $x$};
\draw[-{Latex}] (0,-4.5) -- (0,4.5) node[above] {\small $y$};
\begin{scope}
\clip (-4.3,-4.3) rectangle (-2.08,4.3);
\draw[teal, thick, domain=-4.2:-2.08, samples=60, smooth] plot (\x, {3/((\x+2)*(\x-1))});
\end{scope}
\begin{scope}
\clip (-1.92,-4.3) rectangle (0.92,4.3);
\draw[teal, thick, domain=-1.92:0.92, samples=80, smooth] plot (\x, {3/((\x+2)*(\x-1))});
\end{scope}
\begin{scope}
\clip (1.08,-4.3) rectangle (4.3,4.3);
\draw[teal, thick, domain=1.08:4.2, samples=60, smooth] plot (\x, {3/((\x+2)*(\x-1))});
\end{scope}
\draw[pinkanno, thick] (0,-1.5) circle (2.2pt);
\node[below] at (-2,-4.6) {\small $x=-2$};
\node[below] at (1,-4.6) {\small $x=1$};
\end{scope}
\end{tikzpicture}
\end{center}
\end{friendlyex}

\begin{solutionbox}
Simplify first: $f(x)=\dfrac{3x}{x(x+2)(x-1)}=\dfrac{3}{(x+2)(x-1)}$ for $x\neq 0$.

\textbf{Vertical asymptotes:} $x=-2$ and $x=1$ (from the remaining denominator factors).

\textbf{Hole:} at $x=0$ (the factor $x$ cancels). The hole's $y$-value: $\dfrac{3}{(0+2)(0-1)}=\dfrac{3}{-2}=-1.5$, so the hole is at $(0,-1.5)$ --- notice this means the graph has \textbf{no} $y$-intercept, even though $x=0$ looks like it should give one!

\textbf{Horizontal asymptote:} numerator degree $0<$ denominator degree $2$ (after simplifying), so $y=0$.
\end{solutionbox}

\begin{friendlyex}{}
The graph of $f(x)=\dfrac{2x^2}{x^2-4x+3}$ is shown below. Identify the vertical asymptotes, the horizontal asymptote, and the intercepts.

\begin{center}
\begin{tikzpicture}
\begin{scope}[xscale=0.85,yscale=0.4]
\draw[dashed, gray] (1,-8.5) -- (1,8.5);
\draw[dashed, gray] (3,-8.5) -- (3,8.5);
\draw[dashed, gray] (-4.3,2) -- (5.3,2);
\draw[-{Latex}] (-4.3,0) -- (5.3,0) node[right] {\small $x$};
\draw[-{Latex}] (0,-8.5) -- (0,8.5) node[above] {\small $y$};
\begin{scope}
\clip (-4.3,-8.3) rectangle (0.92,8.3);
\draw[teal, thick, domain=-4.2:0.92, samples=70, smooth] plot (\x, {2*\x*\x/(\x*\x-4*\x+3)});
\end{scope}
\begin{scope}
\clip (1.08,-8.3) rectangle (2.92,8.3);
\draw[teal, thick, domain=1.08:2.92, samples=80, smooth] plot (\x, {2*\x*\x/(\x*\x-4*\x+3)});
\end{scope}
\begin{scope}
\clip (3.08,-8.3) rectangle (5.3,8.3);
\draw[teal, thick, domain=3.08:5.2, samples=60, smooth] plot (\x, {2*\x*\x/(\x*\x-4*\x+3)});
\end{scope}
\filldraw[pinkanno] (0,0) circle (2.2pt);
\node[below] at (1,-8.6) {\small $x=1$};
\node[below] at (3,-8.6) {\small $x=3$};
\node[right] at (5.3,2) {\small $y=2$};
\end{scope}
\end{tikzpicture}
\end{center}
\end{friendlyex}

\begin{solutionbox}
No common factors between $2x^2$ and $x^2-4x+3=(x-1)(x-3)$, so there are no holes.

\textbf{Vertical asymptotes:} $x=1$ and $x=3$.

\textbf{Horizontal asymptote:} numerator and denominator both degree $2$, ratio of leading coefficients $\dfrac{2}{1}=2$, so $y=2$.

\textbf{Intercepts:} $f(0)=\dfrac{0}{3}=0$, so $(0,0)$ is both the $x$- and $y$-intercept. Since $2x^2$ has a double zero at $x=0$, the graph \textbf{touches} the $x$-axis at the origin rather than crossing it.
\end{solutionbox}

\section*{5. Crossing the Horizontal Asymptote}

\begin{friendlyrule}{A Common Misconception}
A graph can \textbf{never} cross a vertical asymptote, but it \textbf{can} cross a horizontal asymptote. To find where, set $f(x)$ equal to the horizontal asymptote's $y$-value and solve.
\end{friendlyrule}

\begin{friendlyex}{}
Find the horizontal asymptote of $f(x)=\dfrac{x+3}{x^2+1}$, and determine where the graph crosses it.
\end{friendlyex}

\begin{solutionbox}
Numerator degree $1<$ denominator degree $2$, so the horizontal asymptote is $y=0$.

Set $f(x)=0$:
\[
\frac{x+3}{x^2+1}=0 \implies x+3=0 \implies x=-3.
\]

The graph crosses its horizontal asymptote at $\boxed{(-3,0)}$.
\end{solutionbox}

\begin{friendlyex}{}
Find the horizontal asymptote of $g(x)=\dfrac{-2x^2+x-4}{x^2+1}$, and determine where the graph crosses it.
\end{friendlyex}

\begin{solutionbox}
Numerator and denominator both degree $2$, ratio of leading coefficients $\dfrac{-2}{1}=-2$, so the horizontal asymptote is $y=-2$.

Set $g(x)=-2$:
\[
\frac{-2x^2+x-4}{x^2+1}=-2 \implies -2x^2+x-4=-2(x^2+1)=-2x^2-2.
\]

The $-2x^2$ terms cancel:
\[
x-4=-2 \implies x=2.
\]

The graph crosses its horizontal asymptote at $\boxed{(2,-2)}$.
\end{solutionbox}

\section*{Summary}

\begin{friendlyrule}{Key Ideas}
\begin{itemize}
    \item The domain of a rational function excludes every $x$ that makes the denominator $0$.
    \item Factor the numerator and denominator completely: a factor that \textbf{cancels} produces a \textbf{hole}; a factor that \textbf{remains} in the denominator produces a \textbf{vertical asymptote}.
    \item Horizontal asymptotes depend on the degrees of the numerator ($n$) and denominator ($m$): $n<m \Rightarrow y=0$; $n=m \Rightarrow y=$ ratio of leading coefficients; $n>m \Rightarrow$ no horizontal asymptote.
    \item A graph can \textbf{never} cross a vertical asymptote, but it \textbf{can} cross a horizontal asymptote --- set $f(x)$ equal to the asymptote's $y$-value and solve.
\end{itemize}
\end{friendlyrule}

\end{document}
